\documentclass[12pt]{article}
\usepackage[margin=1in]{geometry}
\usepackage{enumitem}
\usepackage{titlesec}
\usepackage{parskip}
\usepackage{graphicx}
\usepackage{xcolor}
\usepackage{hyperref}
\usepackage{fontenc}

\title{\textbf{Cheibub et al 2004 RG Notes}\\
\large Cheibub, José Antonio, Adam Przeworski, and Sebastian M. Saiegh. ``Government Coalitions and Legislative Success Under Presidentialism and Parliamentarism'' [2004].\\
\textit{British Journal of Political Science} 34(4): 565-587.}
\author{}
\date{}

\begin{document}

\maketitle

Broad question: can coalitions form under presidential systems? Finds broadly yes, even when the president maintains a minority governing position

Central differentiating factor with presidentialism is that a hostile legislative majority cannot oust the executive: the president's ability to stay in power, at least through the current term, is generally equivalent regardless of whether their current governing party comprises a legislative majority versus minority. (Critique/extension here: what about impeachment or similar proceedings? Seems potentially feasible that there exists some meaningful threshold, although it is presumably a supermajority)

However, the legislative threat may be even greater: if the president is leading via minority, they may essentially lose much policy-making capacity as the legislature remains hostile (it may be electorally beneficial for them to be outright hostile regardless of specific policy preferences)

An important distinction: does the president have agenda-setting authority (i.e., the ability to propose legislation, such as in many Latin American democracies) or largely second-mover responsibilities (i.e., US presidential veto power?) This also has implications for what the status quo reversion might look like (think Romer-Rosenthal)

These guys really hate Linz

\end{document}
